\chapter{Spatial Vector Algebra}
\label{chap:2}
Spatial vectors are 6D vectors that combine the linear and angular aspects of rigid-body motions and forces. They provide a compact notation for studying rigid-body dynamics, in which a single spatial vector can do the work of two 3D vectors, and a single spatial equation replaces two (or sometimes more) 3D vector equations. Spatial vector notation allows us to develop equations of motion quickly, and to express them succinctly in symbolic form. It also allows us to derive dynamics algorithms quickly, and to express them in a compact form that is easily converted into efficient computer code. This chapter presents the spatial vector algebra from first principles up to the equation of motion for a single rigid body. It also presents the planar vector algebra, which is an adaptation of spatial vector algebra to rigid bodies that move only in a 2D plane. Methods of implementing spatial vector arithmetic are covered in Appendix A; and the use of spatial vectors to analyse general rigid-body systems is covered in Chapter 3. 2.1 Mathematical Preliminaries This section reviews briefly some of the mathematical concepts and notations used in spatial vector algebra. More details can be found in books like Greenwood (1988), Selig (1996) and many others. Vectors and Vector Spaces: Linear algebra defines a vector as being an element of a vector space; and different kinds of vector are elements of different vector spaces. Four vector spaces occur frequently in this book, and have been given special names: Rn — coordinate vectors En — Euclidean vectors Mn — spatial motion vectors Fn — spatial force vectors

In each case, the superscript indicates the dimension. One more space that occurs frequently is the space of m × n matrices, which is denoted Rm×n. A coordinate vector is an n-tuple of real numbers, or, in matrix form, an n×1 matrix of real numbers (i.e., a column vector). Coordinate vectors typically represent other vectors; and we use the term abstract vector to refer to the vector being represented. Euclidean vectors have the special property that a Euclidean inner product is defined on them. This product endows them with the familiar properties of magnitude and direction. The 3D vectors used to describe rigid-body dynamics are Euclidean vectors. Spatial vectors are not Euclidean, but are instead the elements of a pair of vector spaces: one for motion vectors and one for forces. Spatial motion vectors describe attributes of rigid-body motion, such as velocity and acceleration, while spatial force vectors describe force, impulse and momentum. The two spaces M6 and F6 are the main topic of this chapter. Hats and Underlines: When spatial vectors and 3D vectors appear together, we mark some spatial vectors with hats (e.g. ˆv, ˆf) in order to avoid name clashes with 3D vectors having the same name. Likewise, we normally make no distinction between coordinate vectors and the abstract vectors they represent, but if a distinction is required then we underline the coordinate vector (e.g. v representing v). We use these notational devices only where they are needed, and hardly ever outside this chapter. The Dual of a Vector Space: Let V be a vector space. Its dual, denoted V ∗, is a vector space having the same dimension as V , and having the property that a scalar product is defined between it and V (i.e., the scalar product takes one argument from each space). If u ∈V ∗and v ∈V then this scalar product can be written either u · v or v · u, the two expressions meaning the same. Duality is a symmetrical relationship: if U = V ∗then V = U ∗. The notion of duality is relevant to spatial vector algebra because the spaces Mn and Fn are dual (i.e., each is the dual of the other). In particular, a scalar product is defined between motion vectors and force vectors such that if m ∈M6 describes the velocity of a rigid body and f ∈F6 describes the force acting on it, then m · f is the power delivered by the force. The scalar product between a vector space and its dual is required to be nondegenerate (also called nonsingular). A scalar product is nondegenerate if it has the following property: for any v ∈V , if v̸ = 0 then there exists at least one vector u ∈V ∗satisfying v · u̸ = 0. This property is a sufficient condition to guarantee the existence of a dual basis on V and V ∗. Dual Bases: Suppose we have two vector spaces, U and V , such that U = V ∗. Let D = \{d1, . . . , dn\} be a basis on U, and let E = \{e1, . . . , en\} be a basis on V . If these bases satisfy the following condition, then they form a dual basis

on U and V : di · ej =  1 if i = j 0 otherwise. This is known as the reciprocity condition. If D and E satisfy this condition, then each is the reciprocal of the other. If one basis is given, then the other is determined uniquely by the reciprocity condition. We can use the notation D∗to denote the reciprocal of D. Note that a dual basis consists of two sets of basis vectors, and covers two vector spaces. If U = En then it becomes possible to set V = U, so that D and E are bases on the same vector space. It also becomes possible to find special cases of D with the property D = E. Such bases are orthonormal. Thus, an orthonormal basis is a special case of a dual basis in which U = V = En and D = E. Dual Coordinates: A dual basis defines a dual coordinate system. We will always use dual coordinates for representing spatial vectors; and we will mostly use Pl¨ucker coordinates, which are a particularly convenient choice of dual coordinate system. The special property of dual coordinates is that u · v = u T v , where u and v are coordinate vectors representing u ∈U and v ∈V in the dual coordinate system defined by D and E. An immediate corollary is that the individual elements in u and v are given by ui = ei · u and vi = di · v . To perform a coordinate transformation from one dual coordinate system to another, we need two transformation matrices: one to operate on u , and one to operate on v . If X is the matrix that performs the desired coordinate transformation on u , then the matrix that does the same job on v is denoted X∗, and the two matrices are related by the equation X∗= X−T . This relationship follows from the requirement that u Tv = (Xu )T(X∗v ) for all u and v . Operators a· and a×: It is possible to interpret expressions like a · b and a × b as being the result of applying an operator, a· or a×, to the operand b. Thus, a· is the operator that maps b to the scalar a · b, while a× maps b to a × b. If a is a coordinate vector, then a· = aT, and a× is a square matrix. Some properties of a× are listed in Tables 2.1 and 2.3. Table 2.3 also lists properties of the operator a×∗, which is the dual of a×.

Dyads and Dyadics: An expression of the form a b· is called a dyad. It is a linear operator that maps vectors to vectors. In particular, the dyad a b· will map the vector c to the vector a (b · c), which is a scalar multiple of a. If a and b are coordinate vectors representing a and b, then the dyad a b· is represented by the matrix a b T, which has a rank of 1. In the most general case, we can have a ∈U, b ∈V and c ∈W, where there are no restrictions on the choice of any of these vector spaces, except that a scalar product must be defined between V and W. Such a dyad maps from W to U. A dyadic (or dyadic tensor) is a general linear mapping of vectors. Any dyadic can be expressed as a sum of r linearly independent dyads: L = r X i=1 ai bi· , where ai ∈U, bi ∈V , and r ≤min(m, n), where m = dim(U) and n = dim(V ). (A set of dyads is linearly independent if the vectors ai are linearly independent in U, and the vectors bi are linearly independent in V .) If ai and bi are coordinate vectors, then L can be written L = r X i=1 ai bT i , which is an m × n matrix of rank r. If m = n = r then L is a 1 : 1 mapping, and is invertible. Dyadics are used in connection with the spatial inertia tensor. 2.2 Spatial Velocity Let us work out a formula for the spatial velocity of a rigid body. We start with a rigid body, B, and choose a fixed point, O, which can be located anywhere in space. (See Figure 2.1(a).) Given O, the velocity of B can be specified by a pair of 3D vectors: the linear velocity, vO, of the body-fixed point that currently coincides with O, and an angular velocity vector, ω. The body as a whole can then be regarded as translating with a linear velocity of vO, while simultaneously rotating with an angular velocity of ω about an axis passing through O. From this description, we can calculate the velocity of any other point in the body using the formula vP = vO + ω × −−→ OP , (2.1) where P is the point of interest, vP is the velocity of the body-fixed point currently at P, and −−→ OP gives the position of P relative to O. This equation shows clearly that there are two contributions to vP : one due to the translation of the whole body by vO, and one due to its rotation by ω about an axis passing

x y z O vO ω B (a) O dOx dOy dOz dx dy dz (b) Figure 2.1: The velocity of a rigid body expressed in terms of ω and vO (a), and the basis vectors for Pl¨ucker motion coordinates (b) through O. Although each individual contribution depends on the position of O, their sum does not (i.e., the dependencies cancel out), so that vP depends only on the motion of B and the location of P. Note that P, vP and ω together describe the same rigid-body motion as O, vO and ω. We now introduce a Cartesian coordinate frame, Oxyz, with its origin at O. This frame defines three mutually perpendicular directions, x, y and z, and three directed lines, Ox, Oy and Oz, each passing through the point O. This frame also defines an orthonormal basis, \{i, j, k\} ⊂E3, which we can use to express ω and vO in terms of their Cartesian coordinates: ω = ωxi + ωyj + ωzk and vO = vOxi + vOyj + vOzk . Having resolved ω and vO into their coordinates, we can describe the velocity of B as the sum of six elementary motions: a rotation of magnitude ωx about the line Ox, a rotation of magnitude ωy about Oy, a rotation of magnitude ωz about Oz, and translations of magnitudes vOx, vOy and vOz in the x, y and z directions, respectively. Our objective is to obtain a spatial velocity vector, ˆv ∈M6, that describes the same motion as ω and vO. We can accomplish this by first defining a basis on M6, and then working out the coordinates of ˆv in that basis. The basis we will use is DO = \{dOx, dOy, dOz, dx, dy, dz\} ⊂M6 , (2.2) which is illustrated in Figure 2.1(b). This is called a Pl¨ucker basis, and it defines a Pl¨ucker coordinate system on M6. The first three basis vectors, dOx, dOy and dOz, are unit rotations about the lines Ox, Oy and Oz, respectively, and the rest are unit translations in the x, y and z directions, respectively. On comparing the definitions of these basis vectors with the six elementary motions just described, it follows immediately that ˆv = ωxdOx + ωydOy + ωzdOz + vOxdx + vOydy + vOzdz . (2.3)

Thus, the Pl¨ucker coordinates of ˆv in DO are the Cartesian coordinates of ω and vO in \{i, j, k\}. Each individual term on the right-hand side depends on the position and orientation of Oxyz, but it can be shown that the expression as a whole is invariant. (See Example 2.2.) The coordinate vector that represents ˆv in DO coordinates can be written ˆv O =

ωx ωy ωz vOx vOy vOz

=  ω v O  , (2.4) where the expression on the far right is a convenient compact notation in which the 6D coordinate vector is expressed as the concatenation of the two 3D coordinate vectors ω = [ ωx ωy ωz ]T and v O = [ vOx vOy vOz ]T. Example 2.1 Equation 2.1 associates a vector vP with each point P in space. It therefore defines a vector field—the velocity field of body B. We could have written this equation in the form V (P) = vO + ω × −−→ OP , (2.5) where V denotes the vector field, and V (P) is the vector associated by V with the point P. Although O appears in the definition of V , the right-hand side is invariant with respect to the location of O, and so the field itself is independent of O. This connection between spatial vectors and vector fields can help one to visualize what a spatial vector really is, and what a spatial-vector expression really means. For example, suppose ˆv1 and ˆv2 are two spatial velocities, and V 1 and V 2 are the corresponding vector fields. If ˆvsum = ˆv1 + ˆv2 then the vector field that corresponds to ˆvsum has the property Vsum(P) = V 1(P) + V 2(P) for all P. Example 2.2 Let us investigate the invariance of Eq. 2.3 with respect to the choice of origin. We can do this by defining ˆv′ to be the spatial velocity obtained by using the point P as origin, and showing that ˆv′ = ˆv. The expression for ˆv′ is therefore ˆv′ = ωxdP x + ωydP y + ωzdP z + vP xdx + vP ydy + vP zdz , and the task is to equate this expression with Eq. 2.3. To keep things simple, we shall choose P to have the coordinates (r, 0, 0) relative to Oxyz, as shown in Figure 2.2. In this case, the linear velocity coordinates are given by Eq. 2.1 as

vP x vP y vP z

= v O + ω ×

r 0 0

=

vOx vOy + ωzr vOz −ωyr

;

O dPy dPz −rdy rdz r P dPx Figure 2.2: Expressing dP y and dP z at O and the three basis vectors that depend on P are dP x = dOx , dP y = dOy + rdz , dP z = dOz −rdy . Combining these equations gives ˆv′ = ωxdOx + ωy(dOy + rdz) + ωz(dOz −rdy) + vOxdx + (vOy + ωzr)dy + (vOz −ωyr)dz , which can be seen to be equal to Eq. 2.3. Returning to the three basis vectors, we have dP x = dOx because the two lines Px and Ox are the same. To understand the expression for dP y, imagine a rigid body that is rotating with unit angular velocity about Py, so that its spatial velocity is dP y. The body-fixed point at O will then have a linear velocity of magnitude r in the z direction, as shown in Figure 2.2; and so the body’s velocity can be expressed as the sum of a unit angular velocity about the line Oy and a linear velocity of magnitude r in the z direction (i.e., dOy +rdz).1 Repeating this exercise with Pz yields the given expression for dP z. The analysis presented here is easily extended to the case of a general location for P, and a similar argument can be used to show that Eq. 2.3 is invariant with respect to the orientation of Oxyz. 2.3 Spatial Force Having obtained an expression for spatial velocity, let us now do the same for spatial force. Given an arbitrary point, O, which can be located anywhere in space, the most general force that can act on a rigid body B consists of a linear force, f, acting along a line that passes through O, together with a couple, nO, 1Readers who are puzzled by this are encouraged to invent and solve their own Pl¨ucker coordinate exercises. A degree of practice with these coordinates is helpful.

B O y z f nO x (a) eOx eOy eOz ex ey ez O (b) Figure 2.3: The force acting on a rigid body expressed in terms of f and nO (a), and the basis vectors for Pl¨ucker force coordinates (b) which is equal to the total moment about O. (See Figure 2.3(a).) Once again, we have a representation comprising one arbitrary point and two 3D vectors. f has taken the place of ω, and nO has taken the place of vO. Given O, nO and f, we can calculate the total moment about any other point, P, using the formula nP = nO + f × −−→ OP . (2.6) This equation is the force-vector analogue of Eq. 2.1, and it shares with that equation the property that the right-hand side is independent of the location of O. The three quantities P, nP and f together describe the same applied force as O, nO and f. We now introduce a Cartesian coordinate frame, Oxyz, with its origin at O, and use it to define the orthonormal basis \{i, j, k\} ∈E3 so that nO and f can be written in terms of their Cartesian coordinates: nO = nOxi + nOyj + nOzk and f = fxi + fyj + fzk . Having resolved nO and f into their components, we can describe the force acting on B as the sum of six elementary forces: couples of magnitudes nOx, nOy and nOz in the x, y and z directions, respectively, and linear forces of magnitudes fx, fy and fz acting along the lines Ox, Oy and Oz, respectively. This coordinate frame also defines the following Pl¨ucker basis on F6: EO = \{ex, ey, ez, eOx, eOy, eOz\} ⊂F6 , (2.7) which is illustrated in Figure 2.3(b). The first three elements are unit couples in the x, y and z directions, and the rest are unit linear forces acting along the lines Ox, Oy and Oz. On comparing the definitions of these basis vectors with the six elementary forces just described, it follows immediately that if ˆf ∈F6 is the spatial force vector representing the same force as O, nO and f, then ˆf = nOxex + nOyey + nOzez + fxeOx + fyeOy + fzeOz . (2.8)

Each individual term on the right-hand side of this equation depends on the position and orientation of Oxyz, but the expression as a whole is invariant. Thus, the value of ˆf does not depend at all on Oxyz, but only on the forces acting on B. The coordinate vector representing ˆf in EO coordinates is ˆf O =

nOx nOy nOz fx fy fz

= n O f  , (2.9) where n O = [nOx nOy nOz]T and f = [fx fy fz]T. 2.4 Pl¨ucker Notation Pl¨ucker coordinates date back to the 19th century, but the basis vectors are new.2 They are the coordinates of choice for 6D vectors, partly because they lend themselves to efficient computer implementation, but mainly because they are the most convenient and easy to use. The notations appearing in Eqs. 2.4 and 2.9 combine Pl¨ucker coordinates with notations that are specific to velocity and force. In the case of velocity, we have used two different symbols (v and ω) to represent linear and angular velocity in 3D vector notation, and we have preserved these symbols in the list of Pl¨ucker coordinates. The story is the same for forces: we have used different symbols for the 3D linear force and moment vectors, and preserved their names in the list of Pl¨ucker coordinates. In general case, the coordinates would have the same name as the vector. Thus, if ˆm and ˆf denote generic elements of M6 and F6, respectively, then their Pl¨ucker coordinates would be ˆm O =

mx my mz mOx mOy mOz

=  m m O  and ˆf O =

fOx fOy fOz fx fy fz

=  f f O  . We shall always list Pl¨ucker coordinates in angular-before-linear order; that is, the three angular coordinates will always be listed first. However, readers should be aware that it is perfectly acceptable to list the three linear coordinates first, and they will encounter this alternative ordering in the literature. From a 2The coordinates appeared in Pl¨ucker (1866). The basis vectors appeared in the author’s Ph.D. thesis, and again in Featherstone (1987), under the name ‘standard basis’. The notation used here is different, but the concept is the same.

mathematical point of view, it makes no difference which order the coordinates are listed. However, from a computational point of view, it does matter because the software is typically written to work with one particular order. 2.5 Line Vectors and Free Vectors A line vector is a quantity that is characterized by a directed line and a magnitude. A pure rotation of a rigid body is a line vector, and so is a linear force acting on a rigid body. A free vector is a quantity that can be characterized by a magnitude and a direction. Pure translations of a rigid body are free vectors, and so are pure couples. A line vector can be specified by five numbers, and a free vector by three. A line vector can also be specified by a free vector and any one point on the line. Let ˆs be any spatial vector, motion or force, and let s and sO be the two 3D coordinate vectors that supply the Pl¨ucker coordinates of ˆs. Here are some basic facts about line vectors and free vectors. • If s = 0 then ˆs is a free vector. • If s · sO = 0 then ˆs is a line vector. The direction of the line is given by s, and the line itself is the set of points P that satisfy −−→ OP × s = sO. • Any spatial vector can be expressed as the sum of a line vector and a free vector. If the line vector must pass through a given point, then the expression is unique. • Any spatial vector, other than a free vector, can be expressed uniquely as the sum of a line vector and a parallel free vector. The expression (for a motion vector) is  s sO −hs  +  0 hs  where h = s · sO s · s . (2.10) This last result implies that any spatial vector, other than a free vector, can be described uniquely by a directed line, a linear magnitude and an angular magnitude. Free vectors can also be described in this manner, but the description is not unique, as only the direction of the line matters. According to screw theory, the most general motion of a rigid body consists of a rotation about a line in space together with a translation along it (i.e., a helical or screwing motion). Such a motion is called a twist. Likewise, the most general force that can act on a rigid body consists of a linear force acting along a line in space, together with a couple acting about it (so the force and couple are parallel). Such a force is called a wrench. In both cases, the line itself is called the screw axis, and the ratio of the free-vector magnitude to the linevector magnitude is called the pitch. Pure free vectors are treated as twists and wrenches of infinite pitch, and do not have a definite screw axis. The elements of M6 and F6 can be regarded as twists and wrenches, respectively.

Magnitude Spatial vectors are not Euclidean, and therefore do not have magnitudes in the Euclidean sense. However, there are four sets of spatial vectors for which a meaningful magnitude can be defined: rotations, translations, linear forces and couples (i.e., all the line vectors and free vectors in M6 and F6). Magnitudes are only comparable within each set; so the magnitude of a rotation, for example, can be compared with that of another rotation, but not with a translation or a force. A vector having unit magnitude can be called a unit vector. However, the meaning of ‘unit magnitude’ does depend on the chosen system of units. For example, a unit rotational velocity is one radian per chosen time unit, and a unit translational velocity is one length unit per time unit. Thus, the Pl¨ucker basis vectors, which are all unit vectors, depend not only on the position and velocity of the coordinate frame, but also on the chosen system of units. 2.6 Scalar Product A scalar product is defined on spatial vectors, in which one argument must be a motion vector, the other must be a force vector, and the result is a scalar representing energy, power or some similar quantity. Given m ∈M6 and f ∈F6, the scalar product can be written either m·f or f ·m, both meaning the same, but the expressions m·m and f·f are not defined (i.e., there is no inner product on either M6 or F6). If f is the force acting on a rigid body, and m is the velocity of that body, then m · f is the power delivered by f. The scalar product creates a connection between M6 and F6 which is formally known as a duality relationship: each space is the dual of the other. The practical consequences of this property for spatial vector algebra are as follows: 1. we use dual coordinates on M6 and F6; 2. force and motion vectors obey different coordinate transformation rules; and 3. there are two cross-product operators: one for motion vectors and one for forces (see §2.9). A dual coordinate system on M6 and F6 is any coordinate system formed by two bases, \{d1, . . . , d6\} ⊂M6 and \{e1, . . . , e6\} ⊂F6, in which the basis vectors satisfy the condition di · ej =  1 if i = j 0 otherwise. (2.11) From the definitions of DO and EO in Eqs. 2.2 and 2.7, it can be seen that these bases form a dual coordinate system. (Treat dOx as d1, dOy as d2, and so on.)

The most important property of a dual coordinate system is that if m and f represent m ∈M6 and f ∈F6 in a dual coordinate system, then m · f = m Tf . (2.12) An immediate consequence of this property is that we need different coordinate transformation matrices for motions and forces. If the matrix X performs a coordinate transformation on motion vectors, and X∗performs the same transformation on force vectors, then the two are related by X∗= X−T . (2.13) This equation follows from the requirement that m Tf = (Xm )T(X∗f ) for all m and f . 2.7 Using Spatial Vectors Spatial vectors are a tool for expressing and analysing the dynamics of rigid bodies and rigid-body systems. This section presents a list of rules for using spatial vectors, which describe the connections between spatial vectors and the physical phenomena they represent. The last few entries preview some rules for quantities that we have not yet met: acceleration, momentum and inertia. The list is not complete. usage: The elements of M6 describe aspects of rigid-body motion, such as: velocity, acceleration, infinitesimal displacement, and directions of motion freedom and constraint. The elements of F6 describe the forces acting on a rigid body, and quantities related to force such as: momentum, impulse, and directions of force freedom and constraint. uniqueness: There is a 1 : 1 mapping between the elements of M6 and the set of all possible motions of a rigid body. Likewise, there is a 1 : 1 mapping between the elements of F6 and the set of all possible forces acting on a rigid body (i.e., the set of all wrenches). relative velocity: If bodies B1 and B2 have velocities v1 and v2, respectively, then the relative velocity of B2 with respect to B1 is vrel = v2 −v1. rigid connection: If two bodies are connected together rigidly, then their velocities are the same. summation of forces: If forces f1 and f2 both act on the same rigid body, then they are equivalent to a single force, ftot, given by ftot = f1 + f2. action and reaction: If body B1 exerts a force of f on body B2, then B2 exerts a force of −f on B1 (Newton’s 3rd law).

scalar product: If a force f acts on a rigid body having a velocity v, then the power delivered by the force is f · v. scaling: A force of αf acting on a body with velocity βv delivers αβ times as much power as a force of f acting on a body with velocity v. A body moving at a constant velocity of v makes the same movement in α units of time as a body moving with velocity αv in one unit of time. acceleration: Spatial acceleration is the rate of change of spatial velocity. Spatial accelerations are true vectors, and follow the same summation rule as velocities. For example, if vrel = v2 −v1 then arel = a2 −a1. summation of inertias: If bodies B1 and B2 have inertias of I1 and I2, and they are connected together to form a single composite rigid body, then the inertia of the composite is I1 + I2. momentum: If a rigid body has a spatial inertia of I and a velocity of v, then its momentum is Iv. equation of motion: The rate of change of a body’s momentum equals the force acting on it; that is, f = d/dt(Iv). If we perform the differentiation, then we get the formula f = Ia + v ×∗Iv. joint 1 s1 s2 sN body 1 body 2 base s3 body N Figure 2.4: Kinematic chain Example 2.3 Figure 2.4 shows a collection of rigid bodies joined together in a chain, with one end joined to a fixed base. A system like this is typical of an industrial robot arm, and is called a kinematic chain. There are a total of N moving bodies and N joints in this system, and the bodies are connected together such that joint i connects from body i −1 to body i. (The base is body 0.) We say that a joint connects ‘from’ one body ‘to’ another so that we can define the velocity across the joint as being the velocity of the ‘to’ body relative to the ‘from’ body. Thus, if vJi is the velocity across joint i, and vi is the velocity of body i, then vJi = vi −vi−1 . (2.14)

We shall assume that each joint allows only a single degree of motion freedom, so that vJi can be described as a scalar multiple of a single motion vector. Specifically, vJi = si ˙qi , (2.15) where si and ˙qi are the axis vector and velocity variable for joint i. If joint i is revolute, then si will be a unit rotation vector so that vJi will be a unit rotational velocity when ˙qi = 1. Likewise, if joint i is prismatic then si will be a unit translation vector. Let us work out the velocity of each body in the system. From Eqs. 2.14 and 2.15 we have vi = vi−1 + si ˙qi . (v0 = 0) (2.16) This equation lets us calculate the velocity of each body in turn, starting at the base and working out towards the free end of the chain. Another expression for the velocity is vi = i X j=1 sj ˙qj . (2.17) From a computational point of view, Eq. 2.16 is more efficient than Eq. 2.17 because it reuses the results of previous calculations. Yet another expression for vi is vi = s1 s2 · · · si 0 · · · 0

˙q1 ... ˙qN

= Ji ˙q . (2.18) In this equation, Ji is a 6 × N matrix called the body Jacobian for body i, and ˙q is the joint-space velocity vector. 2.8 Coordinate Transforms The coordinate transform from A to B coordinates for a motion vector is usually written BXA, and the same transform for a force vector is BX∗ A. The two are related by the equation BX∗ A = BX−T A (cf. Eq. 2.13). In this section, we shall develop the formulae for coordinate transforms between Pl¨ucker coordinate systems. Let A and B denote two Pl¨ucker coordinate systems. Each is defined by the position and orientation of a Cartesian frame, and each frame also defines a Cartesian coordinate system. For convenience, we shall use the same names for all three entities; so frame A defines both Pl¨ucker coordinate system A and Cartesian coordinate system A, and similarly for B. The formula for BXA depends only on the position and orientation of frame B relative to frame A, so it can be expressed as the product of two simpler transforms: one depending on the relative position of frame B and the other on its relative orientation.

Rotation Let A and B be two Cartesian frames with a common origin at O. Given O, any motion vector ˆm ∈M6 can be represented by two 3D vectors, m and mO, and the Pl¨ucker coordinates of ˆm will be the Cartesian coordinates of m and mO. Therefore, let A ˆm, Am, AmO, B ˆm, Bm and BmO be the coordinate vectors representing ˆm, m and mO in A and B coordinates, respectively, and let E be the 3 × 3 rotation matrix that transforms 3D vectors from A to B coordinates (i.e., E = BEA). We now have B ˆm =  Bm BmO  =  E Am E AmO  = E 0 0 E  A ˆm , so BXA = E 0 0 E  . (2.19) The equivalent transform for a force vector is exactly the same: BX∗ A = BX−T A = E 0 0 E  . (2.20) Translation Let O and P be two points in space, and let there be a Cartesian frame located at each point, the two frames having the same orientation. Any motion vector ˆm ∈M6 can be expressed at O by the 3D vectors m and mO, and at P by the vectors m and mP , where mP = mO −−−→ OP ×m (as per Eq. 2.1). In both cases, the Pl¨ucker coordinates of ˆm are the Cartesian coordinates of the appropriate pair of 3D vectors, so we have ˆmP =  m mP  =  m mO −r × m  =  1 0 −r× 1  ˆmO , where r = −−→ OP . Therefore PXO =  1 0 −r× 1  . (2.21) The equivalent transform for a force vector is PX∗ O = PX−T O =  1 −r× 0 1  . (2.22) An expression of the form a× (pronounced ‘a-cross’) denotes an operator that maps b to a × b. If a is a 3D coordinate vector, then a× is the 3 × 3 matrix given by the formula

x y z

× =

0 −z y z 0 −x −y x 0

. (2.23)

a× = −a×T (λ a)× = λ (a×) (λ is a scalar) (a + b)× = a× + b× (E a)× = E a× E−1 (E is orthonormal) (a×) b = −(b×) a (a × b)· = a· b× ((a × b)T = aTb×) (a × b)× = a× b× −b× a× (a × b)× = b aT −a bT a× b× = b aT −(a · b) 13×3 Table 2.1: Properties of the Euclidean cross operator This operator has several mathematical properties, which are listed in Table 2.1. We shall interpret expressions of the form a × b × c as meaning (a×)(b×)c, which is the same as a×(b×c). Any other interpretation requires parentheses. Likewise, the expressions a + b × c and E a× mean a + (b × c) and E(a×), respectively, and any other interpretation requires parentheses. General Transforms Let A and B be Cartesian frames with origins at O and P, respectively; let r be the coordinate vector expressing −−→ OP in A coordinates; and let E be the rotation matrix that transforms 3D vectors from A to B coordinates. The Pl¨ucker transform from A to B coordinates is then the product of a translation by r and a rotation by E. The formula for a motion-vector transform is BXA = E 0 0 E   1 0 −r× 1  =  E 0 −E r× E  ; (2.24) the equivalent transform for a force vector is BX∗ A =  E 0 0 E   1 −r× 0 1  =  E −E r× 0 E  ; (2.25) and the inverses of these two transforms are AXB =  1 0 r× 1  ET 0 0 ET  =  ET 0 r× ET ET  (2.26) and AX∗ B = 1 r× 0 1  ET 0 0 ET  = ET r× ET 0 ET  . (2.27) For comparison, the formula for the 4 × 4 homogeneous coordinate transformation matrix from A to B (homogeneous) coordinates is BTA = E −Er 0 1  . (2.28)

rx(θ) =

1 0 0 0 c s 0 −s c

ry(θ) =

c 0 −s 0 1 0 s 0 c

rz(θ) =

c s 0 −s c 0 0 0 1

( c = cos(θ) , s = sin(θ) ) rot(E) = E 0 0 E  xlt(r) =  1 0 −r× 1  rotx(θ) = rot(rx(θ)) roty(θ) = rot(ry(θ)) rotz(θ) = rot(rz(θ)) Table 2.2: Functions for rotations and Pl¨ucker transforms To facilitate the description of Pl¨ucker transforms, we will occasionally use the functions listed in Table 2.2. In terms of these functions, Eq. 2.24 could have been written BXA = rot(E) xlt(r), and so on. Finally, note that we have adopted the following convention for describing coordinate transforms: the formula for BXA (or for BEA) is associated with a description of how you would have to move frame A so as to make it coincide with frame B. For example, the expression ‘rotx(θ)’ reads like an abbreviated version of the statement ‘rotate about the x axis by an amount θ’, and its value is the coordinate transform BXA for the case where frame B is rotated relative to frame A by an amount θ about their common x axis. Readers should bear in mind that some authors use a different convention. 2.9 Spatial Cross Products Let r ∈E3 be a Euclidean vector that is rotating with an angular velocity of ω, but is not otherwise changing. The derivative of this vector is given by the wellknown formula ˙r = ω × r. In this context, ω× is acting like a differentiation operator: it is mapping r to ˙r. It turns out that this idea can be extended to spatial vectors, but we need to have two operators: one for motions and one for forces. Therefore, let ˆm ∈M6 and ˆf ∈F6 be two spatial vectors that are moving with a velocity of ˆv ∈M6 (e.g. because they are fixed in a rigid body having that velocity), but are not otherwise changing. We now define two new operators, × and ×∗, which are required to satisfy the following equations: ˙ˆm = ˆv × ˆm (2.29) and ˙ˆf = ˆv ×∗ˆf . (2.30) The operator ×∗can be regarded as the dual of ×. The relationship between these two operators is similar to that between the coordinate transform matrices X and X∗.

dOx dOy(t) dOy(t+δt) θ=δt (a) dOx dOy(t) dOy(t+δt) δt (b) × dOx dOy dOz dx dy dz dOx 0 dOz −dOy 0 dz −dy dOy −dOz 0 dOx −dz 0 dx dOz dOy −dOx 0 dy −dx 0 dx 0 dz −dy 0 0 0 dy −dz 0 dx 0 0 0 dz dy −dx 0 0 0 0 ×∗ ex ey ez eOx eOy eOz dOx 0 ez −ey 0 eOz −eOy dOy −ez 0 ex −eOz 0 eOx dOz ey −ex 0 eOy −eOx 0 dx 0 0 0 0 ez −ey dy 0 0 0 −ez 0 ex dz 0 0 0 ey −ex 0 Figure 2.5: Spatial cross-product tables These operators can be defined explicitly by means of the multiplication tables shown in Figure 2.5. Each entry in these tables is the time derivative of the basis vector listed on the top row, assuming it has a velocity equal to the basis vector listed in the left-hand column. For example, Figure 2.5(a) shows a coordinate frame that is moving with a velocity equal to dOx, which means that it is rotating with unit angular velocity about the line Ox. Over a period of δt time units, this motion has no effect on the line Ox, but it causes Oy to rotate by δt radians in the y–z plane. So we have dOx(t + δt) = dOx(t), but dOy(t + δt) = dOy(t) + δt dOz. The time derivatives are therefore d dt dOx = lim δt→0 dOx(t + δt) −dOx(t) δt = lim δt→0 dOx(t) −dOx(t) δt = 0 and d dt dOy = lim δt→0 dOy(t + δt) −dOy(t) δt = lim δt→0 dOy(t) + δt dOz −dOy(t) δt = dOz . This accounts for the first two entries on the first line in the first table. Similarly, Figure 2.5(b) shows a coordinate frame that is moving with a velocity equal to dx, which means that it is translating with unit linear velocity in the x direction. Over a period of δt time units, this motion has no effect on the line Ox, but it causes Oy to shift by δt length units in the x direction. We therefore have

v×∗= −v×T (λ v)× = λ (v×) (λ v)×∗= λ (v×∗) (λ is a scalar) (u + v)× = u× + v× (u + v)×∗= u×∗+ v×∗ (Xv)× = X v× X−1 (Xv)×∗= X∗v×∗(X∗)−1 (Pl¨ucker transform) u × v = −v × u (u × v)· = −v· u×∗ (u ×∗f)· = −f · u× (u × v)× = u× v× −v× u× (u × v)×∗= u×∗v×∗−v×∗u×∗ Table 2.3: Properties of the spatial cross operators dOx(t + δt) = dOx(t), implying that ˙dOx = 0, but dOy(t + δt) = dOy(t) + δt dz, implying that ˙dOy = dz. This accounts for the first two entries on the fourth line in the first table. All other entries are calculated in a similar manner. Using the tables in Figure 2.5, we can deduce that the 6 × 6 matrices representing the operators ˆv× and ˆv×∗in Pl¨ucker coordinates are ˆvO× =  ω vO  × =  ω× 0 vO× ω×  (2.31) and ˆvO×∗=  ω vO  ×∗= ω× vO× 0 ω×  = −(ˆvO×)T , (2.32) and that the spatial cross products of two coordinate vectors are  ω vO  ×  m mO  =  ω × m ω × mO + vO × m  (2.33) and  ω vO  ×∗ fO f  = ω × fO + vO × f ω × f  . (2.34) Various properties of these operators are listed in Table 2.3. 2.10 Differentiation Vectors are differentiated in the same way as scalars. If v(x) is a differentiable vector function of a scalar variable x, then its derivative is given by the formula d dxv(x) = lim δx→0 v(x + δx) −v(x) δx . (2.35) This formula applies to all vectors. In particular, it applies both to abstract vectors and to coordinate vectors. The derivative of a coordinate vector is therefore

d dx(u · v) = du dx · v + u · dv dx d dx(u·) = ( du dx )· d dx(u × v) = du dx × v + u × dv dx d dx(u×) = ( du dx )× d dx(u v·) = du dx v· + u ( dv dx)· d dx(λ v) = dλ dx v + λ dv dx (λ scalar) d dx(L v) = dL dx v + L dv dx (L matrix or dyadic) Table 2.4: Differentiation formulae for vector expressions its componentwise derivative, as can be seen from the following equation: d dx

v1(x) ... vn(x)

= lim δx→0 1 δx

v1(x + δx) −v1(x) ... vn(x + δx) −vn(x)

=

dv1 dx... dvn dx

. The formula also applies to matrices and dyadics. We will be concerned mostly with time derivatives, which we will usually write using dot notation; that is, dv/dt = ˙v, d ˙v/dt = ¨v, and so on. The derivative of a vector is a vector of the same kind; so the derivative of a motion vector is a motion vector, the derivative of a force vector is a force vector, and so on. Differentiation formulae for various vector expressions are shown in Table 2.4. Differentiation in Moving Coordinates The problem of differentiation in moving coordinates can be stated as follows: given the coordinate vector that represents a vector v, find the coordinate vector that represents ˙v. If the basis vectors are constants, then the solution is simply the componentwise derivative of the coordinate vector; otherwise, a formula is required that takes into account the motion of the basis vectors. Let A be the Pl¨ucker coordinate system associated with a Cartesian frame that is moving with a spatial velocity of vA. Suppose we have two coordinate vectors, Am and Af, which represent the spatial vectors m ∈M6 and f ∈F6, respectively, and we wish to find the coordinate vectors that represent ˙m and ˙f in A coordinates. These vectors are given by the formulae Adm dt  = dAm dt + AvA × Am (2.36) and Adf dt  = dAf dt + AvA ×∗Af . (2.37) The expressions on the left-hand sides denote the coordinate vectors representing ˙m and ˙f in A coordinates, and the expressions dAm/dt and dAf/dt denote

the componentwise derivatives of Am and Af, respectively. AvA is the coordinate vector representing vA in A coordinates. For comparison, the well-known formula for the derivative of a 3D Euclidean vector in rotating coordinates, expressed in the same notation, is Adu dt  = dAu dt + AωA × Au . (2.38) In this equation, u ∈E3 is the vector being differentiated, and ωA is the angular velocity of the coordinate frame. Note: In formulating the dynamics algorithms that appear in subsequent chapters, we generally avoid performing differentiation in moving coordinates by formulating the algorithms in stationary coordinate systems that happen to coincide with the moving ones at the current instant. Dot and Ring Notation It is slightly risky to mix dot notation with coordinate vectors because it is not obvious whether a symbol like A ˙m means A(dm/dt) or d(Am)/dt. Nevertheless, we shall take this risk, and make the following definitions: for any abstract vector v, and any coordinate vector Av representing v in A coordinates, A ˙v = Adv dt  (2.39) and A˚v = dAv dt . (2.40) With this notation, Eqs. 2.36 and 2.37 can be written A ˙m = A˚ m + AvA × Am (2.41) and A ˙f = A˚ f + AvA ×∗Af . (2.42) Furthermore, if the identity of the coordinate system is clear from the context, then we need not specify it explicitly. The equations can then be written ˙m = ˚ m + vA × m (2.43) and ˙f = ˚ f + vA ×∗f . (2.44) The symbols ˚ m and ˚ f can be regarded as the apparent rates of change of m and f as viewed by an observer having a velocity of vA.

Example 2.4 Let A and B be two Cartesian frames having velocities of vA and vB, respectively, and let BXA be the Pl¨ucker coordinate transform from A to B coordinates. We can work out the time derivative of BXA using these four equations, which follow from Eqs. 2.39 to 2.41: B ˙m = BXA A ˙m B ˚ m = d dt (BXA Am) =  d dt BXA  Am + BXA A˚ m A ˙m = A˚ m + AvA × Am B ˙m = B ˚ m + BvB × Bm . The derivation proceeds as follows:  d dt BXA  Am = B ˚ m −BXA A˚ m = B ˙m −BvB × Bm −BXA(A ˙m −AvA × Am) = B ˙m −BXAA ˙m −BvB × BXA Am + BXA AvA × Am = (BXA AvA× −BvB× BXA) Am . As this must be true for all Am, we have d dt BXA = B(vA −vB)× BXA . (2.45) 2.11 Acceleration We define the spatial acceleration of a rigid body to be the time derivative of its spatial velocity. Thus, if a body has a spatial velocity of [ωT vT O]T expressed at O, then its spatial acceleration is ˆaO = d dt  ω vO  = lim δt→0 1 δt  ω(t + δt) −ω(t) vO(t + δt) −vO(t)  =  ˙ω ˙vO  . (2.46) This definition may sound obvious, but it contains a surprise for those who are already expert in the traditional method of describing rigid-body acceleration. Consider the body shown in Figure 2.6, which is rotating at a constant angular velocity about a fixed line in space. This body has a constant spatial velocity, and must therefore have zero spatial acceleration; yet almost every point in the body is travelling in a circular path, and is therefore accelerating. This paradox is easily resolved. First, one must realise that spatial acceleration is a property of the body as a whole, rather than a property of individual body-fixed points. Second, one must realise that vO is not the velocity of one particular bodyfixed point, but a measure of the flow of points through O. Thus, ˙vO is not the acceleration of an individual point, but a measure of the rate of change of flow. If the flow is constant, then ˙vO will be zero even if each individual point is accelerating.

P O PO ω vO = ω PO Figure 2.6: Acceleration of a uniformly rotating body r(t) r(t+δt) ω(t+δt) O O’ δt r(t) . Figure 2.7: Computing ˙vO from ¨r Let us compare spatial acceleration with the traditional method of describing acceleration. To this end, we introduce two more quantities: a second point, O′, which is the body-fixed point that coincides with O at time t, and a vector, r, which gives the position of O′ as a function of time. Thus, ˙r is the velocity of O′ and ¨r is its acceleration. In the standard textbooks on classical mechanics, the acceleration of a rigid body would be specified by the two 3D vectors ˙ω and ¨r; so let us examine how ˙vO differs from ¨r. From Eq. 2.46 we have ˙vO = lim δt→0 vO(t + δt) −vO(t) δt . Now, vO(t) = ˙r(t) because O′ coincides with O at time t, but vO(t + δt) is not equal to ˙r(t+δt) because O′ has moved away from O by an amount δt ˙r(t) (see Figure 2.7), so vO(t + δt) is given instead by vO(t + δt) = ˙r(t + δt) −ω × δt ˙r(t) . Combining these equations gives ˙vO = lim δt→0 ˙r(t + δt) −˙r(t) −ω × δt ˙r(t) δt = ¨r −ω × ˙r , (2.47) which implies that ˆaO =  ˙ω ¨r −ω × ˙r  . (2.48)

For comparison, we now define a new 6D vector, ˆa′ O =  ˙ω ¨r  , (2.49) which we shall call the classical acceleration of a rigid body, expressed at O. The name reflects the fact that this vector is composed of the two 3D vectors that are used in the classical description of rigid-body acceleration. Physically, the classical acceleration is the apparent derivative of spatial velocity in a Pl¨ucker coordinate system having its origin at O′, rather than O, meaning that the coordinate frame has a pure linear velocity of ˙r. Thus, according to Eq. 2.43, ˆaO and ˆa′ O are related by ˆaO = ˆa′ O +  0 ˙r  × ˆvO , (2.50) which is easily verified from Eqs. 2.48 and 2.49. The big advantage of spatial acceleration is that it is a true vector, and it obeys the same combination and coordinate transformation rules as velocity. For example, if ˆv1 and ˆv2 denote the spatial velocities of bodies B1 and B2, respectively, and ˆvrel is the relative velocity of B2 with respect to B1, then we already know that ˆvrel = ˆv2 −ˆv1 . The equivalent formula for spatial accelerations is obtained immediately by differentiating the velocity formula: d dt ˆvrel = d dt ˆv2 −d dt ˆv1 ⇒ ˆarel = ˆa2 −ˆa1 . (2.51) Thus, spatial accelerations are composed simply by adding them together, just like velocities. This compares favourably with the traditional formulae for composing accelerations, which involve the use of Coriolis terms. Example 2.5 Example 2.3 presented several formulae for the velocities of the bodies in a kinematic chain. Let us now obtain the corresponding formulae for their accelerations. First, we define aJi to be the acceleration across joint i and ai to be the acceleration of body i; so aJi = ˙vJi and ai = ˙vi. As spatial accelerations are additive, we immediately have aJi = ai −ai−1 , (2.52) which is the time derivative of Eq. 2.14. Next, the derivative of Eq. 2.15 is aJi = si ¨qi + ˙si ˙qi , (2.53) where ¨qi is the acceleration variable for joint i. Given that si represents a joint motion axis that is fixed in body i (and also in body i −1), we have ˙si = vi × si . (2.54)

Combining these three equations gives ai = ai−1 + si ¨qi + vi × si ˙qi , (2.55) which is the derivative of Eq. 2.16. Next, the derivative of Eq. 2.17 is ai = i X j=1 sj ¨qj + vj × sj ˙qj = i X j=1 sj ¨qj + i X j=1 j−1 X k=1 sk × sj ˙qj ˙qk . (2.56) Finally, the derivative of Eq. 2.18 is ai = Ji ¨q + ˙Ji ˙q , (2.57) where ¨q is the joint-space acceleration vector and ˙Ji = [˙s1 ˙s2 · · · ˙si 0 · · · 0]. 2.12 Momentum Figure 2.8 shows a moving rigid body whose centre of mass coincides with the fixed point C at the current instant. This body has a spatial velocity of ˆv, which is expressed at C by the two 3D vectors ω and vC. If the body has a mass of m, and a rotational inertia of ¯IC about its centre of mass, then its momentum is described by the two 3D vectors h = m vC and hC = ¯IC ω, where h specifies the magnitude and direction of the body’s linear momentum, and hC specifies the body’s intrinsic angular momentum. Linear momentum is actually a line vector (like linear force), and its line of action passes through the body’s centre of mass. The body’s moment of momentum about any given point, O, is the sum of its intrinsic angular momentum and the moment of its linear momentum about the given point: hO = hC + −−→ OC × h , (2.58) O C mass: m ω vC h = m vC OC = c rotational inertia: IC hC = IC ω Figure 2.8: Spatial momentum

which is essentially the same formula as Eq. 2.6. Let ˆh denote the spatial momentum of this rigid body, and let ˆhC and ˆhO be the coordinate vectors that represent ˆh in Pl¨ucker coordinate systems with origins at C and O, respectively. These coordinate vectors are given by ˆhC = hC h  =  ¯IC ω m vC  (2.59) and ˆhO = hO h  = ¯IC ω + −−→ OC × m vC m vC  =  1 −−→ OC× 0 1  ˆhC . (2.60) As spatial momentum vectors are elements of F6, they have the same algebraic properties as other spatial force vectors. In particular, they transform according to the formula Bˆh = BX∗ A Aˆh, and a scalar product is defined between momenta and elements of M6. The kinetic energy of the body in Figure 2.8 is 1 2 ˆh · ˆv. 2.13 Inertia The spatial inertia tensor of a rigid body defines the relationship between its velocity and momentum. Spatial inertia is therefore a mapping from M6 to F6. If a rigid body has a spatial velocity of v and a spatial inertia of I, then its momentum is given by the formula h = I v . (2.61) If h and v are coordinate vectors, then I is a 6 × 6 matrix. Let IO and IC be the matrices representing the spatial inertia of a given rigid body in Pl¨ucker coordinate systems having origins at O and C, respectively, where O is an arbitrary point and C coincides with the body’s centre of mass. From Eq. 2.59 we have hC = ¯IC 0 0 m 1  vC , which implies that IC = ¯IC 0 0 m 1  . (2.62) The inertia matrix will take this special form whenever the origin coincides with the centre of mass. From Eq. 2.60 we have hO =  1 c× 0 1  IC vC = 1 c× 0 1  IC  1 0 c×T 1  vO ,

where c = −−→ OC, which implies that IO = ¯IC + m c× c×T m c× m c×T m 1  . (2.63) This is the general form of the inertia matrix in Pl¨ucker coordinates. The quantity ¯IC + m c× c×T is the rotational inertia of the rigid body about O. As ¯IC is a symmetric matrix, it follows that both IC and IO are symmetric. Furthermore, if m > 0 and ¯IC is positive definite, then both IC and IO are positive definite. Dyadic Representation of Inertia An expression of the form fg·, where f, g ∈F6, is a dyad that maps motion vectors to force vectors. Specifically, it maps a motion vector m ∈M6 to the force vector f(g · m), which is a scalar multiple of f. This dyad can be represented in dual coordinates by a 6×6 matrix of the form f g T, where f and g are the coordinate vectors representing f and g in the chosen coordinate system. If f = g then the dyad is said to be symmetric, and its matrix representation is also symmetric. Spatial inertia is a symmetric dyadic tensor that maps from M6 to F6. It can therefore be expressed as the sum of six symmetric dyads as follows: I = 6 X i=1 gi gi· (gi ∈F6) . (2.64) Furthermore, a body’s spatial inertia is a quantity that is fixed in the body: it does not vary except as a function of the body’s motion. It is therefore possible to express I as a sum of dyads in which the individual vectors gi are all fixed in the body to which they refer. Given this expression, we can calculate the time derivative of a spatial inertia as follows: d dt I = 6 X i=1 ( ˙gi gi· + gi ˙gi·) = 6 X i=1 (v ×∗gi) gi· + 6 X i=1 gi (v ×∗gi)· = v×∗ 6 X i=1 gi gi· − 6 X i=1 gi gi· ! v× = v×∗I −I v× , (2.65) where v is the body’s velocity. In obtaining this result, we have used Eq. 2.30 and formulae from Tables 2.3 and 2.4.

dyad coordinate transform mapping e.g. form formula type Mn 7→Mn v× m f· BXA (· · · ) AXB similarity Fn 7→Fn v×∗ f m· BX∗ A (· · · ) AX∗ B similarity Mn 7→Fn I f f· BX∗ A (· · · ) AXB congruence Fn 7→Mn Φ m m· BXA (· · · ) AX∗ B congruence Table 2.5: Properties of dyadics Another property that can be deduced from Eq. 2.64 is the coordinate transformation rule for spatial inertias. Let AI and BI be the matrices representing a given spatial inertia in A and B coordinates, respectively. We therefore have AI = 6 X i=1 Agi AgT i and BI = 6 X i=1 Bgi BgT i , where Agi and Bgi are the coordinate vectors representing gi in A and B coordinates, respectively. Now, Bgi = BX∗ A Agi, so we have BI = 6 X i=1 BX∗ A Agi AgT i (BX∗ A)T = BX∗ A AI AXB . (2.66) Spatial Dyadics Spatial inertia is a mapping from M6 to F6; but there are three more mappings we can define involving these two spaces, making a total of four kinds of spatial dyadic tensor. Some of their basic properties are listed in Table 2.5. The symbol Φ in the examples column denotes an inverse inertia (i.e., Φ = I −1). The column marked ‘dyad form’ shows the pattern of motion and force vectors in the dyads for each kind of tensor. If we think of a dyadic as mapping from the ‘From’ space to the ‘To’ space, then the first vector in the dyad must be an element of the ‘To’ space, and the second vector must be an element in the space that is dual to the ‘From’ space. The coordinate transformation formulae are each labelled as either a similarity transform or a congruence transform. A similarity transform preserves properties such as rank and eigenvalues, but not symmetry or positive definiteness. A congruence transform preserves properties like rank, symmetry and positive definiteness, but not eigenvalues. Properties of Spatial Inertia We have already covered coordinate transforms (Eq. 2.66), the time derivative (Eq. 2.65), the general form of a spatial inertia in Pl¨ucker coordinates (Eq. 2.63)

and the fact that inertia maps velocity to momentum (Eq. 2.61). Here are a few more properties. • Ten parameters are required to define a spatial rigid-body inertia. More general kinds of inertia, such as the articulated-body inertias in Chapter 7, require up to 21 parameters, which is the maximum number of independent values in a symmetric 6 × 6 matrix. • The inertia of a rigid body that is composed of multiple parts is the sum of the inertias of its parts. • The kinetic energy of rigid body having a spatial velocity of v and a spatial inertia of I is T = 1 2 v · I v . (2.67) 2.14 Equation of Motion The spatial equation of motion for a rigid body states that the net force acting on the body equals its rate of change of momentum: f = d dt (Iv) = Ia + (v×∗I −Iv×) v = Ia + v×∗Iv . (2.68) We will often write the equation of motion in the form of an inhomogeneous linear equation f = Ia + p , (2.69) in which I and p are the coefficients and f and a are the variables. In typical use, the coefficients are assumed to be known, and one or both of the variables are unknown. p is called the bias force, and it is simply the value that f must take in order to produce zero acceleration. If f is the net force acting on the body, then p = v×∗Iv. The f in Eq. 2.68 is always the net force acting on the body, but Eq. 2.69 affords the possibility to transfer some known components of the net force from f to p. For example, suppose the net force, fnet, is the sum of an unknown force, fu, and a known gravitational force, fg. Equation 2.68 would then be fnet = Ia + v×∗Iv ; but we have the option of choosing fu as the left-hand quantity in Eq. 2.69, in which case the equation would be fu = Ia + p , and p would be defined as p = v×∗Iv −fg .

Example 2.6 The traditional 3D equations of motion can be recovered from Eq. 2.68 by expanding the spatial quantities into their 3D components. To keep the equations simple, we place the origin at the centre of mass. Equation 2.68 then expands to nC f  = ¯IC 0 0 m 1   ˙ω ¨c −ω × vC  + ω× vC× 0 ω×  ¯IC 0 0 m 1   ω vC  =  ¯IC ˙ω m ¨c −m ω × vC  +  ω × ¯IC ω + m vC × vC m ω × vC  = ¯IC ˙ω + ω × ¯IC ω m ¨c  . (We have used Eq. 2.48 to express the body’s spatial acceleration in terms of ˙ω, ¨c and vC (= ˙c).) Thus, we have recovered Newton’s equation for the motion of the centre of mass, f = m ¨c , (2.70) and also Euler’s equation for the rotational motion of the body about its centre of mass, nC = ¯IC ˙ω + ω × ¯IC ω . (2.71) 2.15 Inverse Inertia It is sometimes useful to write the equation of motion in the form a = Φf + b , (2.72) where Φ is the body’s inverse inertia and b is its bias acceleration, which is the acceleration that the body would have if f = 0. If this equation is the inverse of Eq. 2.69, then Φ = I−1 and b = −Φ p. The advantage of Eq. 2.72 over Eq. 2.69 is that it can be applied to a constrained rigid body. Inverse inertia is a symmetric dyadic tensor that maps from F6 to M6. Some of its properties are already listed in Table 2.5. Another useful property is that the range of an inverse inertia is the subspace of M6 within which the body it applies to is free to move. This, in turn, implies that the rank of an inverse inertia equals the degree of motion freedom of the body to which it applies. Thus, if the body has fewer than six degrees of freedom, then Φ will be singular, but I will not exist. The formulae for the inverse inertia of an unconstrained rigid body, expressed in Pl¨ucker coordinates at C (the centre of mass) and O (a general point) are ΦC = ¯I −1 C 0 0 1 m1  (2.73) and ΦO = " ¯I −1 C ¯I −1 C c×T c× ¯I −1 C 1 m1 + c× ¯I −1 C c×T \# . (2.74)

These are just the inverses of Eqs. 2.62 and 2.63. 2.16 Planar Vectors If a rigid body is constrained to move parallel to a given plane, then it becomes, in effect, a ‘planar’ rigid body. If every body in a rigid-body system is constrained to move parallel to the same plane, then it is a planar rigid-body system. Many practical rigid-body systems are planar. Spatial vectors can easily describe the motions of planar rigid bodies, since planar motion is a subset of spatial motion. However, it can be advantageous to define a stripped-down version of spatial vector algebra specifically for planar rigid bodies—a planar vector algebra. A planar rigid body has three degrees of motion freedom, so planar vectors are elements of the vector spaces M3 and F3. The easiest way to derive planar vector algebra is to start with spatial vector algebra and restrict it to the x–y plane. Essentially, this means removing the first, second and sixth basis vectors and their corresponding coordinates. Thus: dOx dOy dOz dx dy

⇒ dO dx dy dz ex ey ez eOx eOy

⇒ e eOx eOy eOz ωx ωy ωz vOx vOy

⇒ ω vOx vOy vOz nOx nOy nOz fx fy

⇒ nO fx fy fz Note that we have removed the z subscript from the angular basis vectors and coordinates. The planar versions of the spatial cross operator, coordinate transform and inertia matrix can all be obtained by extracting the appropriate 3×3 submatrix from the 6×6 matrix. For example, the planar cross operator matrix is obtained as follows:  ω vO  × =  ω× 0 vO× ω×  =

0 −ωz ωy 0 0 0 ωz 0 −ωx 0 0 0 −ωy ωx 0 0 0 0 0 −vOz vOy 0 −ωz ωy vOz 0 −vOx ωz 0 −ωx −vOy vOx 0 −ωy ωx 0

, so

ω vOx vOy

× =

0 0 0 vOy 0 −ω −vOx ω 0

(2.75) and

ω vOx vOy

×

m mOx mOy

=

0 vOy m −ω mOy ω mOx −vOx m

. (2.76)

In like manner, we can deduce that the formula for a planar coordinate transformation matrix corresponding to a shift of origin from (0, 0) to (x, y), followed by a rotation of the coordinate frame about its new origin by an angle of θ, is X =

1 0 0 0 c s 0 −s c

1 0 0 −y 1 0 x 0 1

=

1 0 0 sx −cy c s cx + sy −s c

, (2.77) where c = cos(θ) and s = sin(θ). Note that planar vectors inherit the properties v×∗= −(v×)T and X∗= X−T from their spatial-vector counterparts. And finally, the formula for the inertia matrix of a planar rigid body is I =

IC + m (c2 x + c2 y) −m cy m cx −m cy m 0 m cx 0 m

=

IO −m cy m cx −m cy m 0 m cx 0 m

, (2.78) where m is the mass, cx and cy are the coordinates of the centre of mass, and IC and IO are the rotational inertias of the body about its centre of mass and coordinate origin, respectively. A total of four parameters are needed to define a planar rigid-body inertia. 2.17 Further Reading Spatial vectors resemble the twists and wrenches of screw theory, the realnumber motors and the elements of the Lie algebra se(3) and its dual, se∗(3). Screw theory was developed in the 19th century, and predates the modern concept of a vector. The main work on screw theory is Ball (1900), but more modern treatments can be found in Angeles (2003), Hunt (1978) and Selig (1996). Motors come in two varieties: 6D vectors over the field of real numbers and 3D vectors over the ring of dual numbers. The former are described in von Mises (1924a,b) and the latter in Brand (1953). The main difference between spatial vectors and motors is that all motors reside in a single vector space. In this respect, the spatial vectors in Featherstone (1987) closely resemble motor algebra. Formally, a Lie algebra is the tangent space to a Lie group at the identity. The spaces M6 and M3 correspond to the Lie algebras se(3) and se(2), and F6 and F3 correspond to their duals, se∗(3) and se∗(2). se(3) and se(2) are the algebras associated with the special Euclidean groups SE(3) and SE(2), which are the sets of every possible displacement of a rigid body in a 3D or 2D Euclidean space. Lie algebras are covered in Murray et al. (1994) and Selig (1996). Other articles that may be of interest include Jain (1991); Park et al. (1995); Pl¨ucker (1866); Rodriguez et al. (1991); Woo and Freudenstein (1970). There are many more, of course—the literature on 6D vectors of various kinds is quite extensive.
